\documentclass[10pt,a4paper]{amsart}
\usepackage[margin=19mm]{geometry}
\usepackage{amsmath,amssymb,mathtools}
\usepackage[scr=rsfs,cal=euler]{mathalfa}
\usepackage{xfrac}
\usepackage[unicode,hidelinks,bookmarks=false]{hyperref}
\newcommand{\ie}{i.e.\ }
\newcommand{\cf}{cf.\ }
\newcommand{\resp}{respectively\ }
\newcommand{\Subsection}{\subsection}
\providecommand{\nobreakdash}{\nobreak\hbox{-}}
\setlength{\emergencystretch}{2em}
\multlinegap0pt
\usepackage{amssymb}
\usepackage{tikz}
\newtheorem{proposition}{Proposition}
\newtheorem{theorem}{Theorem}
\DeclareMathOperator{\Ker}{Ker}
\DeclareMathOperator{\Ran}{Ran}
\title{Finite-Approximate Solvability of Linear Operator Equations}
\author{Nazim I. Mahmudov}
\date{Source: arXiv:2604.22279v1 (CC BY 4.0)}
\begin{document}
\maketitle

\noindent\emph{Source and licence.} Finite-Approximate Solvability of Linear Operator Equations. Version arXiv:2604.22279v1 (CC BY 4.0), distributed by arXiv under the Creative Commons Attribution 4.0 International licence, http://creativecommons.org/licenses/by/4.0/, as recorded in the arXiv licence metadata for this exact version and shown on the abstract page https://arxiv.org/abs/2604.22279v1. Full licence notice: \texttt{licenses/arxiv-2604.22279-CC-BY-4.0.txt}.\par\smallskip

\noindent\textit{Editorial scope.} Counted unit: the article's theorem theorem (source0.tex lines 156--165). The article's complete original proof of it is delivered with it (source0.tex lines 167--271, one \texttt{proof} environment of 105 lines). No mathematics is written, repaired or re-derived: the statement and the proof are the article's own lines, in the article's own notation, and no retained line is substituted or re-flowed.  Retained uncounted as the source-local context the proof needs: lines 111--117.  Nothing was omitted from any retained span, so the package carries no omission record.  No retained line contains a citation, so the wrapper carries no bibliography and no bibliography entry.  No retained line contains a figure environment, an image-inclusion command or drawing code, so no image is delivered and the package declares no asset dependency.  Editorial formatting only, in addition: an AMS \texttt{amsart} preamble that restates the article's own declarations and macros used by the retained lines, namely the environments \texttt{proposition}, \texttt{theorem} on separate counters, together with the standard packages those lines need.  The retained environments print in this document as Proposition 1, Theorem 1 in order of appearance, whereas the article prints the same items with the numbering of its own full document.  The complete native source of the article as distributed in the arXiv e-print for this exact version is bundled unchanged as \texttt{sources/199080/source0.tex}.
	\begin{proposition} \label{prop:1}
		Let \(H\) and \(U\) be Hilbert spaces, let \(L: U \to H\) be a bounded linear operator, set \(\Gamma := L L^*\), and assume \(\Gamma\) is strictly positive. Let \(\pi: H \to H_0\) be an orthogonal projector onto a finite-dimensional subspace \(H_0 \subset H\). Then for any \(\alpha > 0\), the operator
		\[
		T_\alpha := \alpha(I-\pi) + \Gamma
		\]
		is invertible.
	\end{proposition}

	\begin{theorem}[Finite-Approximate Solvability]
		Let \(H\) and \(U\) be Hilbert spaces, let \(L: U \to H\) be a bounded linear operator, set \(\Gamma = LL^*\), and let \(\pi: H \to H_0\) be an orthogonal projector onto a finite-dimensional subspace \(H_0 \subset H\). For any \(h \in H\) and \(\varepsilon > 0\), there exists \(u \in U\) such that
		\[
		\|Lu - h\| < \varepsilon \quad \text{and} \quad \pi(Lu) = \pi h
		\]
		if and only if
		\[
		\lim_{\alpha \to 0^+} \alpha (\alpha(I-\pi) + \Gamma)^{-1} h = 0.
		\]
	\end{theorem}

	\begin{proof}
		Recall that \(\Gamma = LL^*\) is a bounded, self-adjoint, nonnegative operator on \(H\), and \(\pi: H \to H_0 \subset H\) is an orthogonal projector. For \(\alpha > 0\) the bounded operator \(\alpha(I-\pi) + \Gamma\) is invertible (see Proposition \ref{prop:1}, which requires \(\Ker(\Gamma) \cap \Ran(\pi) = \{0\}\)), so the expression below is meaningful.
		
		Define, for \(\alpha > 0\),
		\[
		z_\alpha := (\alpha(I-\pi)+\Gamma)^{-1} h \in H,
		\qquad
		u_\alpha := L^* z_\alpha \in U,
		\qquad
		y_\alpha := \alpha z_\alpha.
		\]
		A useful algebraic identity is obtained by rearranging the equation
		\[
		(\alpha(I-\pi)+\Gamma) z_\alpha = h
		\]
		to obtain
		\begin{equation}\label{eq:basic-id}
			\Gamma(\alpha(I-\pi)+\Gamma)^{-1} = I - \alpha (I-\pi)(\alpha(I-\pi)+\Gamma)^{-1}.
		\end{equation}
		From \eqref{eq:basic-id} we get the following exact formulas for \(u_\alpha\) and its image under \(L\):
		\[
		L u_\alpha = \Gamma z_\alpha = \Gamma(\alpha(I-\pi)+\Gamma)^{-1} h
		= h - \alpha (I-\pi)(\alpha(I-\pi)+\Gamma)^{-1} h
		= h - (I-\pi) y_\alpha.
		\]
		Also, applying \(\pi\) to \eqref{eq:basic-id} and using \(\pi(I-\pi)=0\) gives
		\[
		\pi\Gamma(\alpha(I-\pi)+\Gamma)^{-1} = \pi,
		\]
		hence
		\[
		\pi(Lu_\alpha)=\pi h \qquad\text{for every }\alpha>0.
		\]
		Moreover,
		\begin{equation}\label{eq:error-form}
			Lu_\alpha - h = - (I-\pi) y_\alpha
			= -\alpha (I-\pi)(\alpha(I-\pi)+\Gamma)^{-1} h,
		\end{equation}
		so
		\[
		\|Lu_\alpha-h\| = \alpha\big\|(I-\pi)(\alpha(I-\pi)+\Gamma)^{-1}h\big\|
		\le \big\| \alpha(\alpha(I-\pi)+\Gamma)^{-1} h\big\|.
		\]
		
		\medskip\noindent
		\textbf{(``If'')} Assume
		\[
		\lim_{\alpha\to0^+}\alpha(\alpha(I-\pi)+\Gamma)^{-1}h = 0.
		\]
		Then by \eqref{eq:error-form} we have \(\|Lu_\alpha-h\|\to0\) as \(\alpha\to0^+\), and we already saw \(\pi(Lu_\alpha)=\pi h\) for every \(\alpha>0\). Therefore for any \(\varepsilon>0\) we may choose \(\alpha>0\) sufficiently small so that \(\|Lu_\alpha-h\|<\varepsilon\) while \(\pi(Lu_\alpha)=\pi h\). This proves the existence of the required \(u\) and hence the ``if'' part.
		
		\medskip\noindent
		\textbf{(``Only if'')} We now present the ``only if'' direction using a direct construction for the test vector.
		
		Assume the finite-approximate solvability property holds (so \(h \in \overline{\mathcal S}\) where \(\mathcal S = \{Lu: u \in U,\ \pi(Lu)=\pi h\}\)). Suppose, for contradiction, that
		\[
		\lim_{\alpha\to0^+} y_\alpha = \lim_{\alpha\to0^+} \alpha(\alpha(I-\pi)+\Gamma)^{-1}h \neq 0.
		\]
		Then there exists a sequence \(\alpha_n \downarrow 0\) and a nonzero vector \(y \in H\) with \(y_{\alpha_n} \to y\) in \(H\) and \(y \ne 0\).
		
		First note that \(\pi y_{\alpha_n} \to 0\). Indeed, apply \(\pi\) to the equality
		\[
		(\alpha_n(I-\pi)+\Gamma) z_{\alpha_n} = h.
		\]
		Because \(\pi(I-\pi)=0\) we get
		\[
		\pi\Gamma z_{\alpha_n} = \pi h \qquad\text{for every }n.
		\]
		Restricting \(\Gamma\) to the finite-dimensional space \(\Ran \pi\) shows \(\pi z_{\alpha_n}\) is uniformly bounded in \(n\) (the operator \(\pi\Gamma\pi\) acts on a finite-dimensional space, so the equation \(\pi\Gamma\pi(\pi z_{\alpha_n}) = \pi h - \pi\Gamma(I-\pi)z_{\alpha_n}\) gives uniform boundedness). Hence \(\pi y_{\alpha_n} = \alpha_n \pi z_{\alpha_n} \to 0\) because \(\alpha_n \to 0\) and \(\{\pi z_{\alpha_n}\}\) is bounded. Passing to the limit yields \(\pi y = 0\). Thus the limit vector satisfies
		\[
		(I-\pi)y = y \neq 0,
		\]
		i.e., \(y\) lies in \(\Ker \pi\) (the \((I-\pi)\)-subspace).
		
		Now choose the test vector explicitly as
		\[
		v := (I-\pi)y = y.
		\]
		This is nonzero because \((I-\pi)y = y \neq 0\). From \eqref{eq:error-form} we have, for each \(n\),
		\[
		\langle Lu_{\alpha_n}-h, v\rangle = -\langle (I-\pi)y_{\alpha_n}, v\rangle = -\langle y_{\alpha_n}, (I-\pi)v\rangle.
		\]
		But \((I-\pi)v = (I-\pi)^2 y = (I-\pi)y = y = v\), so
		\[
		\langle Lu_{\alpha_n}-h, v\rangle = -\langle y_{\alpha_n}, v\rangle.
		\]
		Letting \(n \to \infty\) gives
		\[
		\lim_{n\to\infty}\langle Lu_{\alpha_n}-h, v\rangle = -\langle y, v\rangle = -\|y\|^2 \neq 0.
		\]
		Hence there exists \(\delta > 0\) and \(N\) such that for all \(n \ge N\),
		\[
		\big|\langle Lu_{\alpha_n}-h, v\rangle\big| \ge \delta.
		\]
		
		On the other hand, since \(h \in \overline{\mathcal S}\), for any \(\varepsilon > 0\) there is some \(s \in \mathcal S\) with \(\|s-h\| < \varepsilon\). Taking \(\varepsilon\) sufficiently small makes \(|\langle s-h, v\rangle| < \delta/2\). But all \(Lu_{\alpha_n}\) lie in \(\mathcal S\) (because \(\pi(Lu_{\alpha_n}) = \pi h\)), and because \(h\) is in the closure of \(\mathcal S\) there must be members of \(\mathcal S\) (in particular some \(Lu_{\alpha_n}\) for large \(n\)) with \(|\langle Lu_{\alpha_n}-h, v\rangle| < \delta/2\), contradicting the previous lower bound \(|\langle Lu_{\alpha_n}-h, v\rangle| \ge \delta\) for large \(n\).
		
		This contradiction shows that the assumption \(y \ne 0\) is false. Therefore
		\[
		\lim_{\alpha\to0^+}\alpha(\alpha(I-\pi)+\Gamma)^{-1}h = 0,
		\]
		which completes the ``only if'' direction.
		
		Combining with the ``if'' direction yields the equivalence stated in the theorem.
	\end{proof}

\end{document}
